% ============================================================
% Appendix Z.2 — FRFP Case Study (Governance / HR):
% Amazon’s Experimental Recruiting Engine
% ============================================================

\section{FRFP Case Study (Governance / HR): Amazon’s Experimental Recruiting Engine}
\label{app:frfp-amazon-recruiting}

This appendix analyses Amazon’s experimental AI-based recruiting engine
(2014--2017) as an FRFP case study.
The project sought to automate résumé screening and candidate ranking, but
was ultimately abandoned when it was found to systematically disadvantage
female applicants for technical roles.
We examine how the failure can be understood in terms of FRFP’s
explicit/tacit split, boundary design, safe horizons, and institutional
semantics.

\subsection{Background}

Around 2014, Amazon engineers were tasked with building a machine-learning
system to streamline hiring for technical positions.
The idea was that an algorithm would:

\begin{itemize}
  \item ingest résumés from past applicants;
  \item learn patterns associated with candidates who were hired or
        performed well;
  \item score new résumés on a 1--5 star scale and help recruiters prioritise
        candidates.
\end{itemize}

To do this, the team trained models on roughly a decade of historical hiring
data---résumés submitted to Amazon, overwhelmingly from men, reflecting the
male-dominated tech workforce.
Over time, they discovered that the system:

\begin{itemize}
  \item systematically downgraded résumés that included the word ``women’s''
        (e.g.\ ``women’s chess club captain'');
  \item penalised graduates of certain all-women’s colleges;
  \item implicitly treated patterns correlated with male candidates as
        positive signals.
\end{itemize}

Engineers tried to patch obvious markers of gender (e.g.\ manually excluding
features containing ``women’’), but the system continued to find proxy
signals, and confidence in its fairness eroded.
By 2017, Amazon reportedly stopped using the tool for live candidate
screening and eventually shut the project down.

\subsection{Explicit and Tacit Spaces in the Recruiting Engine}

\subsubsection*{Explicit space \(E_{\mathrm{HR}}\)}

In FRFP terms, the engine operated entirely in explicit space:

\begin{itemize}
  \item input artefacts: résumés encoded as text and features; job
        descriptions; historical hiring decisions;
  \item learned model parameters and feature weights;
  \item output artefacts: numeric scores (1--5 stars) for each candidate,
        possibly accompanied by internal diagnostics;
  \item dashboards or interfaces used by recruiters to view rankings.
\end{itemize}

All of this lives in \(E_{\mathrm{HR}}\): artefacts that can be stored,
copied, and processed by machines.

\subsubsection*{Tacit space \(T_{\mathrm{HR}}\)}

Tacit space comprises the epistemic states of:

\begin{itemize}
  \item recruiters and hiring managers:
        their sense of role fit, career trajectories, culture, and potential;
  \item HR leadership and legal/compliance teams:
        their understanding of anti-discrimination law, fairness norms,
        and reputational risk;
  \item organisational leadership:
        their preferences over growth, diversity, and risk trade-offs.
\end{itemize}

These tacit states shape how people interpret explicit signals (résumés,
scores, interviews) and ultimately decide whom to hire.

\subsection{Intended Boundary vs.\ Actual Boundary}

\subsubsection*{Stated intention: AI as pre-screening support}

Public accounts suggest that the recruiting engine was intended to act as:

\begin{itemize}
  \item an internal tool that \emph{ranks} candidates;
  \item a way to speed up high-volume screening;
  \item a decision-support system, not an official replacement for human
        hiring decisions.
\end{itemize}

In FRFP language, the intended boundary was:

\begin{itemize}
  \item AI produces scored candidate lists in explicit space;
  \item humans (recruiters, hiring managers) retain tacit authority to
        select candidates, review résumés, and decide offers.
\end{itemize}

\subsubsection*{De facto boundary: algorithm as gatekeeper}

However, in any scaled recruiting workflow, there is a strong pressure to
treat ranked outputs as a gate:

\begin{itemize}
  \item Recruiters focus on top-ranked candidates first; low-scoring
        candidates may never be reviewed.
  \item The star score becomes the primary triage mechanism for who gets
        interviews.
  \item Over time, recruiters may internalise the model’s patterns and
        align their tacit judgments to it, especially if performance metrics
        reward throughput.
\end{itemize}

Thus, the \emph{effective} boundary tended toward:

\begin{itemize}
  \item \emph{model scores as de facto accept/reject filter},
  \item with tacit human oversight reduced to rare overrides.
\end{itemize}

From an FRFP perspective, this is a classic case of boundary erosion:
a system nominally described as ``support'' functions as a gatekeeper,
pulling correctness authority into the model stack.

\subsection{Protocol-Level Failures in FRFP Terms}

\subsubsection*{Z.2.1 Explicit-Space Bias from Historical Data}

\paragraph{Problem.}
The engine was trained on historical Amazon hiring data:

\begin{itemize}
  \item Most past résumés and hires for technical roles were male.
  \item The model learned that patterns associated with male candidates
        were predictive of positive outcomes.
  \item It penalised explicitly gendered tokens (``women’s chess club'')
        and graduates of certain women’s colleges, even after developers
        tried to remove gender markers.
\end{itemize}

In FRFP terms, the explicit mapping
\[
  f_{\mathrm{recruit}} : \text{résumé features} \to \text{score}
\]
embeds a biased semantics: it effectively treats ``being similar to past,
mostly male hires'' as a proxy for fitness.

\paragraph{Consequence.}
Even if humans at the boundary were perfectly vigilant, the explicit artefacts
(scores) already compress candidates into a shape that bakes in historical
bias.
The FRFP kernel would classify this as a defect in explicit semantics:
the pipeline in \(N \ltimes E_{\mathrm{HR}}\) does not implement a
normatively acceptable mapping from résumés to evaluative signals.

\subsubsection*{Z.2.2 Boundary Design and Over-Delegation}

\paragraph{Problem.}
The project’s ambition was to automate as much screening as possible.
In practice, this implies:

\begin{itemize}
  \item using the model to filter or prioritise thousands of résumés;
  \item trusting that top-ranked candidates are ``good enough'' to proceed;
  \item allowing low-ranked candidates to be silently dropped.
\end{itemize}

Under FRFP, this is a high-risk boundary design:

\begin{itemize}
  \item the boundary between explicit scores and tacit acceptance is
        effectively moved upstream into the model;
  \item human recruiters are placed in a role where they mostly confirm
        model suggestions, not critically re-evaluate them;
  \item there is no explicit requirement that humans systematically examine
        candidates the model dislikes to detect systemic bias.
\end{itemize}

\paragraph{Consequence.}
Humans are nominally in the loop, but their tacit authority is weakened.
The model’s preferences dominate which résumés get attention at all,
and thus which tacit evaluations can even occur.
FRFP would characterise this as \emph{boundary capture} by the explicit
pipeline.

\subsubsection*{Z.2.3 Absence of Safe Horizons and Re-Grounding}

\paragraph{Problem.}
Recruiting is inherently long-horizon:

\begin{itemize}
  \item thousands of decisions per year;
  \item gradual evolution of candidate pools, job markets, and company needs;
  \item long-term effects on workforce composition and culture.
\end{itemize}

The engine was conceived as a permanent, automated component of the hiring
pipeline:
once deployed, it would continuously rank incoming résumés with no explicit
protocol for:

\begin{itemize}
  \item periodic re-grounding in normative standards
        (e.g.\ fairness, diversity goals, legal constraints);
  \item systematic audit of its treatment of protected groups;
  \item resetting or recalibrating tacit recruiter expectations relative to
        human baselines.
\end{itemize}

No explicit safe-horizon design appears in public accounts.

\paragraph{Consequence.}
According to FRFP’s dynamics, such a system is structurally prone to:

\begin{itemize}
  \item tacit drift: recruiters gradually internalise the model’s biased
        notion of ``strong candidates'';
  \item institutional entrenchment: over time, it becomes harder to imagine
        hiring without the tool, because everything from SLAs to headcount
        is tuned assuming its throughput;
  \item invisibility of error: individual biased decisions are hard to detect,
        but their cumulative effect on workforce composition is large.
\end{itemize}

In practice, the project was terminated once engineers identified the bias
and concluded it could not easily be fixed.
From an FRFP viewpoint, this termination is an implicit acknowledgement that
safe horizons had already been violated at the design level.

\subsubsection*{Z.2.4 Institutional Governance as Explicit-Only}

\paragraph{Problem.}
The project was an internal experiment.
There is no public evidence of:

\begin{itemize}
  \item a formal, FRFP-style governance body for AI in hiring that defined
        admissible uses, boundary conditions, or safe horizons;
  \item prospective fairness and bias evaluation protocols that had to be
        passed before deployment;
  \item ongoing institutional operators that combined tacit HR/legal judgment
        with model metrics to decide on continued use.
\end{itemize}

Instead, governance appears to have been:

\begin{itemize}
  \item largely internal to the engineering group;
  \item driven by technical performance and ad hoc bias checks;
  \item ultimately resolved by a management decision to scrap the tool
        once bias was recognised.
\end{itemize}

\paragraph{Consequence.}
This matches FRFP’s picture of \emph{explicit-only governance}:
decisions are based on model metrics and internal analyses, without a
structurally distinct tacit governance layer with clear authority and
constraints.
Such governance is precisely what FRFP’s institutional non-automation results
warn about.

\subsection{Relation to FRFP’s Population-Level Concerns}

\subsubsection*{Bias Propagation in Populations}

Subsequent work on AI-assisted hiring suggests that:

\begin{itemize}
  \item humans tend to follow biased model recommendations at high rates
        (e.g.\ $\sim 90\%$ acceptance of biased suggestions in controlled
        experiments);
  \item bias can thus propagate through populations of decision-makers when
        AI systems are used at scale;
  \item humans alone, without structural constraints, are poor correctors of
        embedded algorithmic bias.
\end{itemize}

These findings resonate with FRFP’s population-level semantics:

\begin{itemize}
  \item tacit states of many agents are updated by interacting with biased
        explicit pipelines;
  \item institutional operators that rely on human ``oversight'' without
        explicit boundary and safe-horizon design are insufficient;
  \item biased explicit artefacts can reshape what counts as normal,
        even for agents who initially recognise the bias.
\end{itemize}

\subsubsection*{Non-Automation of Fair Governance}

The Amazon case illustrates a broader claim:

\begin{quote}
  There is no purely explicit procedure---no amount of historical data or
  scoring metrics---that can, by itself, guarantee fair and normatively
  acceptable recruitment decisions over long horizons.
\end{quote}

FRFP frames this as a non-automation result:

\begin{itemize}
  \item fairness, diversity, and legal compliance are normative properties
        of decisions as they interact with broader social structures;
  \item these properties cannot be fully read off from historical hiring data
        or model outputs;
  \item institutional tacit judgment (HR, legal, leadership) must remain
        essential, and must be structurally embedded at the boundary.
\end{itemize}

In the Amazon case, the lack of such an institutional FRFP layer meant that
bias was discovered late and addressed bluntly (by scrapping the tool),
rather than being continuously managed.

\subsection{Counterfactual: An FRFP-Aligned Recruiting Engine}

We now sketch what a recruiting engine might look like under FRFP constraints.

\subsubsection*{Z.2* Explicit/Tacit Roles}

\begin{itemize}
  \item The engine is explicitly defined as a \emph{signal generator}:
        it extracts features, clusters candidates, and proposes structured
        summaries of experience, but does not emit a single canonical
        ``fitness score''.
  \item Tacit evaluation of candidate suitability remains with recruiters and
        hiring managers; their judgments are treated as boundary events and
        logged.
\end{itemize}

\subsubsection*{Z.2* Boundary Design}

\begin{itemize}
  \item No candidate is rejected \emph{solely} on the basis of an AI score;
        low-ranked candidates are either:
        \begin{itemize}
          \item sampled for human review; or
          \item subjected to explicit fairness audits.
        \end{itemize}
  \item UI surfaces AI judgments as \emph{one feature among many}, alongside
        raw résumé content and structured summaries.
  \item At decision time (shortlist, reject, offer), humans must explicitly
        confirm or override AI-influenced decisions, with a brief rationale.
\end{itemize}

\subsubsection*{Z.2* Safe Horizons and Audits}

\begin{itemize}
  \item Safe horizons cap how many decisions can be made under a fixed model
        and dataset before mandatory fairness audits and re-training.
  \item Regular cohort-based audits check disparities across gender and other
        protected attributes in both model scores and final hiring outcomes.
  \item Persistent disparities trigger protocol changes, not just model tweaks:
        e.g.\ additional features, diverse training data, changes in boundary
        rules for low-ranked candidates.
\end{itemize}

\subsubsection*{Z.2* Institutional Governance}

\begin{itemize}
  \item An AI-in-HR governance group defines admissible uses, documents
        explicit/tacit roles, and sets fairness and transparency requirements.
  \item HR, legal, and diversity stakeholders participate in boundary
        design and safe-horizon tuning.
  \item Deployments and updates require governance sign-off, with logs
        and metrics available for internal and external scrutiny.
\end{itemize}

\subsection{Lessons for FRFP and HR / Governance}

The Amazon recruiting engine case supports several FRFP themes:

\begin{enumerate}
  \item \textbf{Historical explicit artefacts encode tacit injustice.}
        Training on past hiring decisions without structural correction
        imports institutional bias into explicit semantics.

  \item \textbf{Without clear boundaries, support tools become gatekeepers.}
        In high-volume contexts like recruiting, ranking and scoring almost
        inevitably take on decision-making roles unless carefully bounded.

  \item \textbf{Safe horizons and fairness audits are structural, not optional.}
        Long-horizon use of biased models can change who gets considered at
        all, reshaping the workforce; FRFP insists on finite horizons and
        periodic tacit re-grounding.

  \item \textbf{Governance cannot be left to the model builders alone.}
        Institutional operators must combine tacit expertise in law, ethics,
        and HR with explicit metrics; otherwise explicit-only governance is
        likely to miss structurally unfair designs.

  \item \textbf{Formal architectures clarify what went wrong.}
        Instead of treating the episode as a one-off ``sexist AI'', FRFP
        reveals a deeper pattern:
        biased explicit training data, boundary capture, absence of safe
        horizons, and weak institutional semantics together made the project
        epistemically and ethically unsound.
\end{enumerate}

As AI-based hiring tools continue to proliferate in industry and on hiring
platforms, this case study illustrates why FRFP-style architectures---with
explicit/tacit separation, clear boundaries, safe horizons, and institutional
oversight---are necessary to avoid repeating the same structural mistakes
at greater scale.
